% Independent continuation, 12 October 2026.
% Fragment uses amsmath, amssymb, amsthm; theorem/lemma/corollary environments.
% Primary bibliography keys below are suggestions for the parent article.
\section{Joint Gamma--polygamma cumulants}
\label{gc:section}

The latest incoming Gamma-normalization report \cite{GammaIncoming} evaluates the appropriate
Gustafson determinant, establishes a paired parameter deformation, and
proves Gaussian moments for a single trigamma statistic through degree
$n+1$. Its research question Q2 asks for a precise classification of
joint polygamma moments. The theorem below answers this question in a
complete weighted range. It also includes the digamma statistic and all
even polygamma orders, which require the unpaired deformation.
The classical integral and the source's density are inherited; the
explicit cumulant classification and its consequences are the present
extension. No claim of priority over all literature is made.

For $n\ge1$, $\sigma>0$, put $N=2n+2$ and consider the probability
measure on $(0,\infty)^n$ with density proportional to
\begin{equation}\label{gc:density}
 \Delta(x_1^2,\ldots,x_n^2)
 \Delta(\cosh\pi x_1,\ldots,\cosh\pi x_n)
 \prod_{j=1}^n x_j\sinh(\pi x_j)
       |\Gamma(\sigma+ix_j/2)|^{2N},
\end{equation}
where $\Delta(y)=\prod_{j<k}(y_k-y_j)$. Its normalizer is
\begin{equation}\label{gc:normalizer}
 Z_{n,\sigma}
 =2^{n(3n+1)/2}\pi^{n(n+1)}n!
       \frac{\Gamma(2\sigma)^{(n+1)(2n+1)}}{\Gamma(N\sigma)}.
\end{equation}
The Vandermondes have the same sign, so this is a positive measure.
Write $\mathbb E=\mathbb E_{n,\sigma}$ and define, for $q\ge0$,
\begin{equation}\label{gc:R}
 R_q=\sum_{j=1}^n\Re\psi^{(q)}(\sigma+ix_j/2),
 \qquad \psi^{(0)}=\psi.
\end{equation}
The \emph{weight} of $R_q$ is $q+1$. A joint cumulant means
\[
 \kappa(Y_1,\ldots,Y_r)
 =\left.\partial_{t_1}\cdots\partial_{t_r}
     \log\mathbb E e^{t_1Y_1+\cdots+t_rY_r}\right|_{\mathbf t=0}.
\]
All fixed moments exist. Indeed the positive-order polygammas are bounded
on each vertical line $\Re z=\sigma$, the digamma grows at most
logarithmically, and the Gamma weight has an exponential margin after
multiplication by every fixed polynomial or logarithmic power.

\begin{theorem}[Explicit joint cumulants in the complete weighted range]
\label{gc:main}
Let $q_1,\ldots,q_r\ge0$ and suppose
\begin{equation}\label{gc:range}
                   \sum_{i=1}^r(q_i+1)\le N=2n+2.
\end{equation}
The one-point and two-point cumulants are
\begin{align}
 \mathbb E R_q
   &=\frac12\left[(N-2^q)\psi^{(q)}(2\sigma)
               -\delta_{q0}\psi(N\sigma)\right],
              \label{gc:mean}\\
 \operatorname{Cov}(R_p,R_q)
   &=\frac14\left[\psi^{(p+q+1)}(2\sigma)
               -\delta_{p0}\delta_{q0}\psi'(N\sigma)\right].
              \label{gc:cov}
\end{align}
For every $r\ge3$ in the same weighted range,
\begin{equation}\label{gc:higher}
 \boxed{\quad
 \kappa(R_{q_1},\ldots,R_{q_r})
 =-\frac{\psi^{(r-1)}(N\sigma)}{2^r}
                    \prod_{i=1}^r\delta_{q_i0}.
 \quad}
\end{equation}
Thus every mixed cumulant of order at least three that contains a
positive-order polygamma statistic vanishes in this range.
\end{theorem}

\subsection{Proof by independent parameter shifts}

For independent parameters $z_1,\ldots,z_N$ near zero, deform all Gamma
parameters to $\sigma+z_\ell$. Gustafson's symplectic Barnes integral,
in the normalization of Derkachov--Manashov \cite[Eq.~(2)]{DerkachovManashov}, and the
source's determinant-to-density transformation give
\begin{equation}\label{gc:unpaired}
\begin{split}
 &\mathbb E\prod_{j=1}^n\prod_{\ell=1}^N
 \frac{\Gamma(\sigma+z_\ell+ix_j/2)
       \Gamma(\sigma+z_\ell-ix_j/2)}
      {\Gamma(\sigma+ix_j/2)\Gamma(\sigma-ix_j/2)}\\
 &\qquad=
 \frac{\Gamma(N\sigma)}{\Gamma(N\sigma+\sum_\ell z_\ell)}
 \prod_{\ell<h}
       \frac{\Gamma(2\sigma+z_\ell+z_h)}{\Gamma(2\sigma)}.
\end{split}
\end{equation}
This is an identity of holomorphic germs. On a sufficiently small
polydisc every shifted Gamma parameter remains in the right half-plane.
Uniform Stirling bounds leave an integrable exponential majorant, even
after any fixed mixed derivative; they justify coefficient extraction
under the integral. Holomorphic log-Gamma is used on that half-plane,
and the analytic logarithm of the expectation is fixed by its value
zero at the origin.

Put $p_k=\sum_{\ell=1}^N z_\ell^k$. Taylor expansion of the logarithm
of the integrand in \eqref{gc:unpaired} gives
\[
 \log\left(\prod_{j,\ell}\frac{\Gamma(\sigma+z_\ell+ix_j/2)
     \Gamma(\sigma+z_\ell-ix_j/2)}
     {\Gamma(\sigma+ix_j/2)\Gamma(\sigma-ix_j/2)}\right)
            =2\sum_{k\ge1}\frac{R_{k-1}}{k!}p_k.
\]
For every integer $k\ge1$ the elementary pair-sum identity is
\begin{equation}\label{gc:pairsum}
 \sum_{\ell<h}(z_\ell+z_h)^k
 =(N-2^{k-1})p_k+\frac12\sum_{a=1}^{k-1}
                              \binom ka p_a p_{k-a}.
\end{equation}
The coefficient of $p_k$ follows by excluding the diagonal from
$\sum_{\ell,h}z_\ell^a z_h^{k-a}$; the two endpoint terms each
contribute $(N-1)p_k$ before the factor $1/2$ is applied.
Taking logarithms in \eqref{gc:unpaired} therefore gives
\begin{equation}\label{gc:power-sum-log}
\begin{split}
 &\log\mathbb E\exp\left(2\sum_{k\ge1}\frac{R_{k-1}}{k!}p_k\right)\\
 &\quad=\sum_{k\ge1}
       \frac{(N-2^{k-1})\psi^{(k-1)}(2\sigma)}{k!}p_k
       +\frac12\sum_{a,b\ge1}
          \frac{\psi^{(a+b-1)}(2\sigma)}{a!b!}p_a p_b
       -\sum_{r\ge1}\frac{\psi^{(r-1)}(N\sigma)}{r!}p_1^r.
\end{split}
\end{equation}

The next finite algebra observation is the step that permits separation
of the different statistics, rather than merely recording a generating
identity.

\begin{lemma}[Independent power-sum coordinates through degree $N$]
\label{gc:psindependent}
Assign weight $k$ to an indeterminate $P_k$. Substitution
$P_k\mapsto p_k(z_1,\ldots,z_N)$ is injective on the space of
polynomials of weighted degree at most $N$.
\end{lemma}
\begin{proof}
Only $P_1,\ldots,P_N$ can occur in that degree range.
Newton's identities give an invertible triangular change between
$p_1,\ldots,p_N$ and the first $N$ elementary symmetric polynomials,
which are algebraically independent. Alternatively the Jacobian of
the power sums is
\[
 \det\left(\frac{\partial p_k}{\partial z_\ell}\right)_{k,\ell=1}^N
              =N!\prod_{\ell<h}(z_h-z_\ell),
\]
which is not the zero polynomial. Either proof gives the claimed
injectivity.
\end{proof}

\begin{proof}[Proof of Theorem~\ref{gc:main}]
The homogeneous component of degree $d\le N$ in
\eqref{gc:power-sum-log} involves only $p_1,\ldots,p_d$.
By Lemma~\ref{gc:psindependent}, compare its coefficients as polynomials
in independent power-sum variables. Set $p_k=k!t_{k-1}/2$ in this
finite-degree comparison. The logarithmic moment generating function
then has, through weighted degree $N$, the form
\begin{equation}\label{gc:cgf}
\begin{split}
 \log\mathbb E\exp\left(\sum_{q\ge0}t_qR_q\right)
  ={}&\frac12\sum_{q\ge0}(N-2^q)\psi^{(q)}(2\sigma)t_q\\
    &+\frac18\sum_{p,q\ge0}\psi^{(p+q+1)}(2\sigma)t_pt_q
      -\sum_{r\ge1}\frac{\psi^{(r-1)}(N\sigma)}{2^r r!}t_0^r.
\end{split}
\end{equation}
Here equality means coefficient equality only through that weighted
degree, not an unrestricted identity of generating functions.
The linear, quadratic and higher coefficients give
\eqref{gc:mean}--\eqref{gc:higher}.
\end{proof}

\subsection{A finite formula for every retained ordinary moment}

Let $\mathcal P_r$ be the set of partitions of $\{1,\ldots,r\}$.
For a block $B$ define $c(B)$ by \eqref{gc:mean} when $|B|=1$,
by \eqref{gc:cov} when $|B|=2$, and by
\eqref{gc:higher} when $|B|\ge3$, using the orders attached to the
labels in $B$. Since every block satisfies \eqref{gc:range} whenever
the full list does, the ordinary moment is explicitly
\begin{equation}\label{gc:moment}
 \boxed{\qquad
 \mathbb E\prod_{i=1}^rR_{q_i}
             =\sum_{\pi\in\mathcal P_r}\prod_{B\in\pi}c(B),
 \qquad \sum_i(q_i+1)\le N.
 \qquad}
\end{equation}
This finite partition formula evaluates all the indicated integrals
using only polygamma values at $2\sigma$ and $N\sigma$. It contains no
unevaluated spectral derivative or auxiliary integral.

For positive orders define the centered statistics
\[
 Y_q=R_q-\frac12(N-2^q)\psi^{(q)}(2\sigma),\qquad q\ge1.
\]
All blocks of size other than two disappear from their centered moments.
Consequently
\begin{corollary}[Weighted Wick rule]
\label{gc:wick}
For $q_i\ge1$ and $\sum_i(q_i+1)\le N$,
\begin{equation}\label{gc:wickeq}
 \mathbb E\prod_{i=1}^rY_{q_i}
 =\begin{cases}
 0,&r\text{ odd},\\[2pt]
 \displaystyle 4^{-r/2}\sum_{\text{pairings }\pi}
         \prod_{\{i,j\}\in\pi}\psi^{(q_i+q_j+1)}(2\sigma),
                                    &r\text{ even}.
 \end{cases}
\end{equation}
\end{corollary}
For example, for $n\ge5$ the genuinely mixed fourth moment is
\begin{equation}\label{gc:example13}
 \mathbb E(Y_1^2Y_3^2)
 =\frac{\psi'''(2\sigma)\psi^{(7)}(2\sigma)}{16}
                +\frac{\psi^{(5)}(2\sigma)^2}{8}.
\end{equation}
At $\sigma=1/4$ this is $33\pi^{12}/2$, using
$\psi'''(1/2)=\pi^4$, $\psi^{(5)}(1/2)=8\pi^6$, and
$\psi^{(7)}(1/2)=136\pi^8$.
For $n\ge2$ another mixed evaluation is
\begin{equation}\label{gc:example012}
 \mathbb E\bigl[(R_0-\mathbb ER_0)(R_1-\mathbb ER_1)
                                  (R_2-\mathbb ER_2)\bigr]=0,
\end{equation}
whose total weight is $1+2+3=6$.
The odd-order restriction, $S_{2a-1}=R_{2a-1}$, reads
\[
 \mathbb ES_{2a-1}=(n+1-4^{a-1})\psi^{(2a-1)}(2\sigma),\qquad
 \operatorname{Cov}(S_{2a-1},S_{2b-1})
               =\tfrac14\psi^{(2a+2b-1)}(2\sigma)
\]
whenever, respectively, $a\le n+1$ and $a+b\le n+1$.
Every higher cumulant vanishes when the sum of the corresponding
$a$-indices is at most $n+1$. This is the requested separation of the
joint odd-polygamma coordinates.

\subsection{New scalar integral evaluations}

Set $n=1$, $\sigma=1/4$ and
\[
 \omega(x)=x\sinh(\pi x)|\Gamma(1/4+ix/2)|^8,
 \quad z_x=1/4+ix/2,
\]
\[
 g(x)=\Re\psi(z_x)+\gamma+3\log2,\qquad
 h(x)=\Re\psi'(z_x)-\pi^2/2.
\]
Here $\int_0^\infty\omega(x)\,dx=4\pi^5$.
The theorem, whose weighted budget is $N=4$, gives the exact ordinary
integrals, using the standard special values in \cite[\S5.15]{DLMF},
\begin{align}
 \int_0^\infty\omega(x)\Re\psi''(z_x)\,dx&=0,
                                                    \label{gc:scalar1}\\
 \int_0^\infty\omega(x)g(x)\Re\psi''(z_x)\,dx&=\pi^9,
                                                    \label{gc:scalar2}\\
 \int_0^\infty\omega(x)g(x)^2h(x)\,dx&=0,
                                                    \label{gc:scalar3}\\
 \int_0^\infty\omega(x)g(x)^3\,dx&=\pi^5\zeta(3),
                                                    \label{gc:scalar4}\\
 \int_0^\infty\omega(x)g(x)^4\,dx&=\frac{\pi^9}{15}.
                                                    \label{gc:scalar5}
\end{align}
For the last two, the centered digamma cumulants are
\[
 \kappa_2(R_0)=\frac{\pi^2}{12},\quad
 \kappa_3(R_0)=\frac{\zeta(3)}4,\quad
 \kappa_4(R_0)=-\frac{\pi^4}{240}.
\]
Thus the fourth central moment is
$3(\pi^2/12)^2-\pi^4/240=\pi^4/60$; multiplication by
$4\pi^5$ gives \eqref{gc:scalar5}.
The nonzero digamma cumulants explain why the Gaussian description
belongs to the positive-order sector.

\subsection{The finite range is essential}

At degree $N+1$, Newton's relation $e_{N+1}=0$ in $N$ variables creates
a nonzero kernel:
\begin{equation}\label{gc:firstrelation}
 \sum_{\lambda\vdash N+1}
       \frac{(-1)^{N+1-\ell(\lambda)}}{z_\lambda}
                 p_{\lambda_1}\cdots p_{\lambda_{\ell(\lambda)}}=0,
 \qquad z_\lambda=\prod_{j\ge1}j^{m_j}m_j!.
\end{equation}
One therefore cannot compare all formal power-sum coefficients beyond
the theorem's range. It is an algebraic failure of this extraction
argument, not a proof that every first excluded cumulant is nonzero.
The following limiting example supplies an actual obstruction to
extending the source's single-trigamma Gaussian law universally.

\begin{proposition}[A nonzero first excluded Gaussian moment]
\label{gc:sharpness}
For $n=1$ define the source-normalized statistic
\[
 X_{1,\sigma}=\frac{2R_1-2\psi'(2\sigma)}
                         {\sqrt{\psi'''(2\sigma)}}.
\]
Then
\begin{equation}\label{gc:skewlimit}
        \lim_{\sigma\downarrow0}\mathbb E X_{1,\sigma}^3
                       =\frac{2\sqrt6}{9}>0.
\end{equation}
In particular the Gaussian matching degree $n+1$ cannot be increased
to $n+2$ uniformly over all $n\ge1$ and all $\sigma>0$.
This proposition does not evaluate the general $(n+2)$nd defect.
\end{proposition}
\begin{proof}
At $n=1$, \eqref{gc:normalizer} is
$Z_{1,\sigma}=4\pi^2\Gamma(2\sigma)^6/\Gamma(4\sigma)$,
and $\sigma^5 Z_{1,\sigma}\to\pi^2/4$.
Set $x=2\sigma y$. Using $z\Gamma(z)\to1$ as $z\to0$, the
normalized density in the $y$ variable tends pointwise to
\[
                q_0(y)=\frac{32}{\pi}\frac{y^2}{(1+y^2)^4},
                \qquad y>0.
\]
Its integral is one by the beta integral. Since the approximating
densities also integrate to one, Scheff\'e's lemma gives convergence
in $L^1$. Also
\[
 \sigma^2\Re\psi'(\sigma(1+iy))
       \longrightarrow\Re(1+iy)^{-2},\qquad
 \sigma^2|\psi'(\sigma(1+iy))|
       \le\sigma^2\psi'(\sigma),
\]
and the upper bound tends to one. The density convergence and this
uniform bound give convergence of every fixed moment of the scaled
statistic. If $t=y^2/(1+y^2)$ then $t$ has the
$\operatorname{Beta}(3/2,5/2)$ law and
\[
 \Re(1+iy)^{-2}=f(t)=1-3t+2t^2.
\]
The beta moment formula
$\mathbb Et^j=(3/2)_j/(4)_j$ gives
\[
 \mathbb Ef=\tfrac14,\qquad
 \mathbb E(f-1/4)^2=\tfrac3{32},\qquad
 \mathbb E(f-1/4)^3=\tfrac1{64}.
\]
Finally $\sigma^2\psi'(2\sigma)\to1/4$ and
$\sigma^4\psi'''(2\sigma)\to3/8$. Thus the normalized variable
converges in moments to $\sqrt{32/3}(f-1/4)$ and its third moment
is $(32/3)^{3/2}/64=2\sqrt6/9$.
\end{proof}

\subsection{A joint Gaussian consequence}

Although the main conclusions are exact integral evaluations, the
finite moment law also has a useful joint limit. Fix distinct positive
orders $q_1,\ldots,q_d$ and define
\[
 C(\sigma)_{ij}=\tfrac14\psi^{(q_i+q_j+1)}(2\sigma).
\]
This matrix is positive definite: the polygamma integral representation
gives
\[
 v^TC(\sigma)v
 =\frac14\int_0^\infty
       \left(\sum_i(-1)^{q_i}v_i t^{q_i}\right)^2
                    \frac{t e^{-2\sigma t}}{1-e^{-t}}\,dt>0
\]
for every nonzero real $v$. Distinct monomial degrees ensure strictness.
For any positive sequence $\sigma_n$, let
\[
 \mathbf Z_n=C(\sigma_n)^{-1/2}
       (R_{q_i}-\mathbb ER_{q_i})_{i=1}^d.
\]
Every mixed moment of total ordinary degree $k$ is exactly the
corresponding standard $d$-variate Gaussian moment once
$k(1+\max_iq_i)\le2n+2$. This follows by expanding the deterministic
linear change of coordinates and applying the weighted Wick rule.
In particular $\mathbf Z_n$ converges in distribution to $N(0,I_d)$.
For a proof, every fixed linear projection has eventually the exact
Gaussian moments of each prescribed order; the second moments give
tightness and higher even moments give uniform integrability.
Moment determinacy of the normal law and the Cram\'er--Wold theorem
finish the argument. No estimate uniform in $\sigma_n$ is needed.

\subsection{Independent verification}

The accompanying program \texttt{verify\_gamma\_cumulants.py} performs
122 exact assertions. These verify the pair-sum identity at twelve
orders and eight numbers of variables, verify the ranks of the
power-sum substitution through the first excluded degree for up to five
variables, and independently evaluate the beta moments and displayed
quarter-parameter arithmetic. In particular the rank drops by one at
weighted degree $N+1$, as predicted by \eqref{gc:firstrelation}.

The numerical side uses continuous Cauchy--Binet in the opposite
direction: it integrates one-dimensional moment entries with inserted
products of polygammas, constructs a determinant of their finite Taylor
jets, and extracts moments and cumulants by finite polynomial
arithmetic. It does not use the Gamma quotient for the expectation
side. Seventeen checks include the zero mean of $R_2$, the third and
fourth centered digamma cumulants, mixed vanishing cumulants, and
$\operatorname{Var}(R_2)=\operatorname{Cov}(R_1,R_3)$ at $n=2$.
The cases use quarter and non-quarter shifts. At 36 decimal digits,
with 56 Gauss--Legendre nodes on each stated integration interval, the
largest scaled residual is approximately $8.2\times10^{-34}$.
A repeat at 46 decimal digits and 72 nodes per interval
reduces the largest scaled residual to approximately $2.1\times10^{-43}$.
A separate direct adaptive integral gives
\[
 \begin{aligned}
 \mathbb EX_{1,0.02}^3
 &=0.5443940332249657482536718308092\ldots,\\
 \mathbb EX_{1,0.005}^3
 &=0.5443321558595570316756456585706\ldots.
 \end{aligned}
\]
consistent with the proved limit
$2\sqrt6/9=0.5443310539518173551549520166013\ldots$.
These floating point residuals are independent representation checks,
not certified error intervals or replacements for the analytic proofs.
